% ============================================================
% §2.3  Open-World Autonomy: Capabilities and Gaps
% ============================================================

\section{Open-World Autonomy: Responses to Mismatch and Their Limits}
\label{sec:capabilities_and_gaps}

We now narrow to the subfield in which this dissertation sits: autonomous agents and robots that must act in open worlds. The central question is practical: when an agent encounters a task-relevant mismatch, what can it do?

We first clarify what ``open world'' means for a deployed agent. We then organize current systems by \emph{how} they respond to mismatch and by the limitation each response encounters under structural novelty. We close by stating the specific gaps this dissertation addresses and the boundaries of its claims. Consistent with the chapter's framing, the discussion is organized around responses to novelty rather than around separate perception, planning, and control pipelines.

\subsection{What ``Open World'' Means for Deployed Agents}
\label{sec:what_is_open_world}

A closed-world agent is one whose model is assumed complete and fixed at design time; an open-world agent is deployed where that assumption fails, encountering objects, dynamics, or task structure not represented when it was built~\citep{langley2020open}. \citet{kejriwal2024challenges} give this intuition a useful structure, distinguishing \emph{weak} open-world learning (detecting that an unexpected change has occurred), \emph{semi-strong} open-world learning (characterizing it), and \emph{strong} open-world learning (adapting to it). They further argue that a mature system should be \emph{antifragile}, improving through encounters with novelty rather than merely surviving them. This detect--characterize--adapt progression maps onto the lifecycle an agent must support; the formal taxonomy adopted in this dissertation is developed in \cref{ch:foundations}.

One terminological caution is worth making explicit: ``open world'' is used in two distinct senses. In this dissertation it means \emph{incompleteness and novelty}: the agent encounters what its models do not cover. Elsewhere it denotes a large, freely explorable \emph{environment}. Platforms such as ODYSSEY, with its library of primitive and composite Minecraft skills~\citep{liu2025odyssey}, and Virtual Community, a multi-agent simulator of human--robot society~\citep{zhou2025virtual}, are ``open worlds'' in this second, spatial sense. The two senses are easily conflated; we use the first throughout.

\subsection{Responses to Mismatch and Their Limits}
\label{sec:response_types}

Despite the diversity of open-world methods, the ways an agent can respond to task-relevant mismatch fall into a small number of types. Each is effective within a range, and each encounters a characteristic limitation under \emph{structural} novelty (\cref{tab:response_types}).

\begin{table}[t]
\centering
\caption{Response types for open-world mismatch and their characteristic limitation under structural novelty (new entities, predicates, or required skills). The last column asks whether the response, on its own, can extend the agent's representation rather than only adjust within it.}
\label{tab:response_types}
\small
\begin{tabular}{@{}p{2.3cm}p{5.4cm}p{2.0cm}@{}}
\toprule
\textbf{Response} & \textbf{Characteristic limitation} & \textbf{Extends repr.?} \\
\midrule
Detect / flag & Signals change; neither characterizes nor repairs it & No \\
Re-estimate & Updates parameters of a fixed model; assumes the vocabulary still fits & No \\
Explore & Seeks new experience; undirected search is slow to localize the fault & Partially \\
Replan / repair & Recombines existing operators; cannot introduce operators it lacks & No \\
End-to-end generalize & Broad priors absorb variation; opaque to failure attribution & Sometimes \\
\bottomrule
\end{tabular}
\end{table}

\emph{Detection} is the entry point: open-set and open-world recognition frameworks flag inputs or categories outside training experience~\citep{boult2021towards}. A perception system may mark an object as unknown, or an execution monitor may notice that an action's observed effect differs from its expected effect. By construction, however, detection only signals that something changed; it neither characterizes the change nor repairs behavior.

\emph{Re-estimation} covers the continual-learning, transfer, and drift methods of \cref{sec:drift_vs_novelty}: it updates a model whose structure is assumed adequate, and so addresses distributional change rather than structural novelty. For example, a controller may adapt to a heavier object or a noisier sensor while retaining the same state and action spaces.

\emph{Exploration}, driven by intrinsic motivation (\cref{sec:intrinsic_motivation}), can in principle discover new behavior, but undirected exploration is expensive and slow to localize which part of the model failed. The method in \cref{ch:rapid_learn} addresses this limitation by using the failed plan's structure to focus exploration.

\emph{Replanning and plan repair} exploit symbolic structure: when execution fails, a planner searches for an alternative arrangement of its existing operators~\citep{fox2006plan, fritz2007monitoring}. This is powerful when a solution exists within the current operator set. If a hallway is blocked, the planner may route around it; if one object is unavailable, it may substitute another object whose modeled effects are equivalent. But plan repair cannot manufacture an operator the agent lacks---exactly what an obstructive novelty demands.

Recent work supplies missing knowledge from outside the planner. Ding et al.'s COWP augments a classical task planner with common-sense knowledge drawn from a language model to handle open-world ``situations,'' evaluated across 561 dinner-serving scenarios in which a normally valid plan becomes inapplicable~\citep{ding2022robot}. Hierarchical formulations instead compose known skills to solve unfamiliar problems, motivated---as \citet{eppe2021hierarchical} observe---by the gap between animals and machines: New Caledonian crows solve metatool problems, keeping out-of-sight tools and subgoals in mind while planning several tool-use steps ahead, whereas RL agents perform well when rules are rigid, as in games, but poorly when small changes in the environment or required actions impair performance.

Finally, \emph{end-to-end methods} rely on broad priors, increasingly from foundation models, to absorb variation without explicit adaptation. Their promise and limits are the subject of \cref{sec:foundation_models}.

A theme cuts across these responses. Approaches that \emph{grow} the representation tend to do so slowly and globally: the DREAM architecture makes this explicit through a redescription cycle that restructures representations on a slower timescale than policy learning~\citep{doncieux2020dream}. Approaches that act \emph{quickly} tend not to grow the representation at all; Balloch et al.'s WorldCloner adapts to injected novelty in a single shot by reusing a learned symbolic transition model~\citep{balloch2023neuro}, but assumes fixed observation and action spaces throughout.

The gap this dissertation targets lies between these poles: fast, \emph{localized} extension of an agent's competence in response to a specific structural failure.

\subsection{Gaps This Dissertation Addresses---and What It Does Not}
\label{sec:ow_gaps}

Three gaps motivate the contributions that follow.

First, the field lacks \emph{systematic evaluation} of structural novelty in sequential decision-making. Open-set and OOD benchmarks measure detection of novel inputs, not how a decision-making agent degrades and recovers when the structure of its task changes. \Cref{ch:benchmarks} addresses this gap with environments that inject controlled structural novelty and metrics for the resulting adaptation dynamics.

Second, the field lacks mechanisms for \emph{failure localization and targeted skill acquisition}: replanning cannot introduce missing operators, and undirected exploration is inefficient. \Cref{ch:rapid_learn} uses the symbolic structure of a failed plan to localize the fault and instantiate a focused learning problem at exactly that point, as in the axe example where the failed \texttt{break} step identifies the behavior that must be repaired.

Third, the field lacks \emph{representations that generalize across object and embodiment variation without full retraining}: much that appears as novelty under a robot-centric representation need not be novel under a better-chosen one. \Cref{ch:generalization} pursues this gap through interaction abstractions that factor learned skills away from particular objects and embodiments.

It is equally important to bound the claims. This dissertation does not contribute open-world \emph{perception}; it assumes access to task-relevant state and treats the perception loop as out of scope, a limitation revisited in \cref{ch:conclusion}. It does not address \emph{safety-critical} continual learning, where transient constraint violations during adaptation are unacceptable~\citep{coursey2026safe}. Nor does it address the \emph{beneficial} exploitation of novelty as opportunity rather than obstacle; the methods developed here are failure-driven, which is an honest limitation.

Finally, it does not address the \emph{normative} dimension of open-world autonomy. Agents deployed among people encounter novel conflicts among goals, norms, and constraints that static policies do not resolve, a problem given dedicated treatment by \citet{jones2026oamncc}, who characterize the knowledge an agent needs to resolve such conflicts in a human-aligned way and frame them as ill-structured problems requiring dynamic reframing. These are real parts of the broader open-world problem. This dissertation addresses a specific, foundational slice: how structured abstractions let an agent adapt to and generalize across structural novelty, rather than how to build a full-stack open-world agent.
